% =====================================================================
% Testing & Performance Evaluation  (~340 words; many numbers are
% placeholders pending full bench-test data)
% =====================================================================
\section{Testing and Performance Evaluation}
\label{sec:testing}

The brief assigns 20\,marks to performance evaluation and asks for
explicit metrics on accuracy, latency, repeatability and robustness.
A four-part test plan was therefore designed against those metrics.
Final numerical results are pending the full bench campaign and will
be inserted into Table~\ref{tab:results} when complete; the test
methods themselves are described below.

\subsection{Functional verification}
A scripted sequence drives each joint independently through its full
travel and back, and then exercises a coordinated multi-joint move
through five waypoints. The test is run in both manual and
autonomous modes and is repeated $N=10$ times. Pass criteria are:
no checksum failures over the run, all waypoints reached within
$\pm 2^\circ$ of the commanded angle, and no link-loss watchdog
trips.

\subsection{Closed-loop accuracy and repeatability}
With a fiducial mounted to the gripper, the arm is commanded to a
fixed target pose 30~times. The mean Euclidean position error at the
end-effector and the standard deviation across trials give the
accuracy and repeatability figures respectively, mapped against
ISO~9283-style metrics~\cite{iso9283}. The same procedure is
repeated with a 50\,g payload to characterise robustness to load
variation; the closed-loop is expected to mask most of the resulting
servo droop.

\subsection{End-to-end latency}
The controller asserts a debug pulse on a spare GPIO when a frame is
emitted; the arm asserts another debug pulse when the corresponding
servo-command call returns. The two pulses are captured on a shared
two-channel oscilloscope and the time delta is recorded. The
expected mean is in the 30--60\,ms range based on
literature~\cite{gomez2012}; a histogram across 200 frames gives the
mean, standard deviation and tail behaviour.

\subsection{Power and link robustness}
Steady-state current draw is measured at the Mega's $V_{in}$ and at
the servo rail under three conditions: idle, manual-mode tracking,
and autonomous sort. Link robustness is tested by sweeping the
controller-to-arm distance up to 10\,m and logging the per-frame
checksum-failure rate.

\begin{table}[t]
  \centering
  \caption{Performance metric summary (placeholders).}
  \label{tab:results}
  \footnotesize
  \begin{tabularx}{\columnwidth}{lXl}
    \toprule
    \textbf{Metric} & \textbf{Method} & \textbf{Result} \\
    \midrule
    Accuracy ($^\circ$)        & per-joint, $N=30$           & \emph{TBC} \\
    Repeatability (1\,$\sigma$)& per-joint, $N=30$           & \emph{TBC} \\
    End-to-end latency (ms)    & GPIO pulse, $N=200$         & \emph{TBC} \\
    Checksum-fail rate (\%)    & 5\,m line of sight, 1\,min  & \emph{TBC} \\
    Idle current (mA)          & Mega + servos               & \emph{TBC} \\
    \bottomrule
  \end{tabularx}
\end{table}

\corehit{Brief Objective~9: evaluate accuracy, latency, repeatability
and robustness; methodology defined and instrumented.}

\advhit{Quantitative bench instrumentation (GPIO pulse pair on a
shared scope; 200-frame latency histograms; payload-disturbance
repeatability under closed loop) used to drive iterative
optimisation rather than ad-hoc demonstration.}
